\section*{Project contract}
Implement \texttt{fir\_q15}. Pass a full input-window address in \texttt{a0}, a full
coefficient address in \texttt{a1}, a full destination address in \texttt{a2} and an
output count \(0\leq n\leq64\) in \texttt{a3}. Return the original count in
\texttt{v0} and the number of rounded results that required saturation in \texttt{v1}.
All stored samples and coefficients are signed sixteen-bit, little endian, at arbitrary
byte alignment.

The input contains seven history samples followed by \(n\) new samples: \(n+7\)
halfwords in chronological order. The address points to the oldest history sample,
not the first new sample. Eight coefficients \(c_0,\ldots,c_7\) are stored in the
same oldest-to-newest window order. Their numerical weights are \(c_j/32768\).
For each output index \(i\), first compute the exact integer accumulator
\[
  S_i=\sum_{j=0}^{7}u_{i+j}c_j,
  \qquad R_i=\left\lfloor\frac{S_i+16384}{32768}\right\rfloor.
\]
Clamp \(R_i\) to \([-32768,32767]\). Count it as clipped only when \(R_i\) lies
strictly outside that interval. Round once after the complete exact sum. The rule
chooses the nearest integer, with halfway cases going toward positive infinity.

Input, coefficient and output slices are pairwise disjoint. Their lengths are
\(2(n+7)\), 16 and \(2n\) bytes. Only logical output bytes may change, including
temporary writes. All caller-data reads fit the selected input, coefficients or output.
Mapped padding does not extend those slices. At zero count, perform no caller-data
read or write, including no coefficient or history read. Private native stack storage
and submission-owned constants are allowed.

Preserve all bits of \texttt{s0} through \texttt{s7}, \texttt{fp} and \texttt{sp},
return through the supplied \texttt{ra}, and produce no console output. Argument,
temporary and vector registers and overlapping floating-point storage may change.
Keep the supplied caller/data harness. No kernel history is retained between calls.
